\documentclass[border=10pt]{standalone}
\usepackage{tikz}
\usetikzlibrary{arrows.meta, positioning, calc}

\begin{document}

\definecolor{sfblue}{RGB}{0, 120, 200}
\definecolor{sfred}{RGB}{220, 50, 30}
\definecolor{sfgreen}{RGB}{40, 160, 40}
\definecolor{sfgray}{RGB}{80, 80, 80}
\definecolor{sflight}{RGB}{220, 240, 255}

\begin{tikzpicture}[
    node distance=1.5cm and 1.5cm,
    block/.style={draw, #1, line width=1.2pt, fill=#1!8,
                  minimum width=2.0cm, minimum height=0.9cm,
                  font=\normalsize\bfseries, rounded corners=3pt},
    block/.default=sfblue,
    gain/.style={draw, #1, line width=1.2pt, fill=#1!8,
                 circle, minimum size=0.9cm,
                 font=\normalsize\bfseries},
    gain/.default=sfblue,
    % Summing-junction circle diameter ~40% of block height (Block Diagram
    % Style convention; matches feedback_block.tex / pid_block.tex).
    % 加え合わせ点の円の直径はブロック高さの約4割（ブロック線図の流儀。
    % feedback_block.tex / pid_block.tex と揃える）。
    sum/.style={draw, sfblue, line width=1.2pt, fill=white,
                circle, minimum size=0.4cm,
                font=\normalsize\bfseries},
    input/.style={font=\normalsize\bfseries, sfgray},
    output/.style={font=\large\bfseries, sfblue},
    line/.style={-{Stealth[length=2.5mm]}, line width=1.2pt, sfgray},
    pathlbl/.style={font=\normalsize\bfseries},
  ]

  % --- Main path (gyro row): omega -> x(Delta t) -> Sigma1 -> alpha -> Sigma2 -> output ---
  % 主経路（ジャイロ行）: ω → ×Δt → Σ1 → α → Σ2 → 出力
  \node[input] (gyro_in) {Gyro $\omega$};
  \node[block=sfred, right=1.3cm of gyro_in] (dt_block) {$\times\Delta t$};
  \node[sum, right=1.3cm of dt_block] (sum1) {$\Sigma$};
  \node[gain=sfred, right=1.3cm of sum1] (alpha) {$\alpha$};
  \node[sum, right=1.3cm of alpha] (sum2) {$\Sigma$};
  \node[output, right=1.1cm of sum2] (output) {$\hat{\theta}$};

  % --- Accel path (top row): drops into Sigma2 from above, entirely outside
  %     the gyro feedback loop below (which stays self-contained around
  %     Sigma1). This is the layout fix versus the previous revision, where
  %     the accel row sat below the main row and visually nested inside the
  %     feedback loop even though it carries no feedback signal. ---
  % 加速度経路（上段）: Σ2 に上から入る。下側のジャイロ・フィードバック
  % ループ（Σ1周りで閉じる）とは完全に別の場所を通る。旧版は加速度経路が
  % 主経路の下段にあり，フィードバック信号ではないのに見た目上ループの
  % 内側に入り込んでいた。これを是正したのが今回のレイアウト。
  \node[input, above=2.0cm of gyro_in] (accel_in) {Accel};
  \node[block=sfgreen, right=1.3cm of accel_in] (atan) {atan2};
  \node[gain=sfgreen, right=1.3cm of atan] (one_alpha) {$1\!-\!\alpha$};

  % --- Feedback delay path (below row): previous estimate theta_hat_{k-1} ---
  % フィードバック遅延経路（下段）: 前回推定値 θ̂_{k-1} を戻す
  \node[block=sfblue, below=1.5cm of sum1, minimum width=1.6cm] (delay) {$z^{-1}$};

  % --- Branch point on the output line (theta_hat is tapped, not re-derived) ---
  % 出力線上の分岐点（θ̂ をそのまま分岐、再計算しない）
  \coordinate (tap) at ($(sum2)!0.5!(output)$);

  % --- Connections: gyro path ---
  \draw[line, sfred] (gyro_in) -- (dt_block);
  \draw[line, sfred] (dt_block) -- (sum1.west);
  \draw[line, sfred] (sum1.east) -- (alpha.west);
  \draw[line, sfred] (alpha.east) -- (sum2.west);

  % --- Connections: accel path. Drops straight down into Sigma2 from the
  %     north; a short right-angle jog (down, then across) absorbs the
  %     small misalignment between one_alpha's x-position and sum2's,
  %     since "right=1.3cm of" spacing is computed from each row's own
  %     node widths and the two rows use different-width nodes. ---
  % 加速度経路の接続。Σ2 に北側からまっすぐ下りる。one_alpha と sum2 の
  % x座標のわずかなズレ（"right=1.3cm of" の間隔は各行のノード幅から
  % 計算されるため，幅の違う行同士では厳密には揃わない）を、
  % 短い直角ジョグ（下りてから横に寄せる）で吸収する。
  \draw[line, sfgreen] (accel_in) -- (atan);
  \draw[line, sfgreen] (atan) -- (one_alpha);
  \draw[line, sfgreen] (one_alpha.south) -- ++(0,-0.4) -| (sum2.north);

  % --- Output line and its branch dot ---
  \draw[line, sfblue, line width=1.5pt] (sum2) -- (output);
  \fill[sfblue] (tap) circle (2.2pt);

  % --- Feedback path: branch dot -> straight down -> into delay from the
  %     east -> up into Sigma1 from below. With the accel path now routed
  %     above the main row, the lower half-plane belongs to this path
  %     alone, so (unlike the previous revision) no detour around another
  %     signal is needed. ---
  % フィードバック経路: 分岐点 → まっすぐ下へ → 遅延ブロックへ東側から
  % 入る → 上へ → Σ1 に南側から入る。加速度経路が上段に移ったことで
  % 下半分はこの経路専用になり，旧版で必要だった他信号を避ける迂回が
  % 不要になった。
  \draw[line, sfblue] (tap) -- (tap |- delay.east) -- (delay.east);
  \draw[line, sfblue] (delay.north) -- (sum1.south);

  % --- Sum node signs (Block Diagram Style convention): "+" sits above a
  %     horizontally (west-)entering arrow. For a vertically-entering arrow
  %     (into Sigma1 from below, into Sigma2 from above) the "+" sits to
  %     the arrow's right, mirroring Sigma1's existing bottom-entry sign. ---
  % 加え合わせ点の符号（ブロック線図の流儀）: 水平（西側）から入る矢印の
  % 「＋」はその上に置く。垂直に入る矢印（Σ1は下から，Σ2は上から）の
  % 「＋」はその右側に置く（Σ1の下からの入力の符号と同じ規則）。
  \node[font=\scriptsize, sfred,  anchor=south east] at ($(sum1.north)+(150:0.3)$) {$+$};
  \node[font=\scriptsize, sfblue, anchor=north west] at ($(sum1.south)+(330:0.3)$) {$+$};
  \node[font=\scriptsize, sfred,  anchor=south east] at ($(sum2.north)+(150:0.3)$) {$+$};
  \node[font=\scriptsize, sfgreen, anchor=south west] at ($(sum2.north)+(30:0.3)$) {$+$};

  % --- Path labels (kept off the signal lines) ---
  \node[pathlbl, sfred, below=0.2cm of dt_block] {High-pass};
  \node[pathlbl, sfgreen, above=0.2cm of atan] {Low-pass};
  \node[pathlbl, sfblue, below=0.15cm of delay] {前回値};

\end{tikzpicture}
\end{document}
